\section{Building Roads}
\label{sec:cses-building-roads}

\problemstatement{Given $n$ cities and $m$ existing roads, find the minimum
number of new roads needed so that every city is reachable from every
other, and output which roads to build.}

\begin{patternbox}
\emph{``Connect all components with the fewest new edges''} reduces to
counting connected components: if there are $c$ components, exactly $c-1$
new roads suffice -- one linking a representative of each component to a
common one.
\end{patternbox}

\subsection*{Components, then bridge them}

Find the connected components (by DFS/BFS or Disjoint Set). Pick one
representative city per component. Building a road from the first
component's representative to each other representative connects everything
with $c-1$ roads, which is optimal -- no fewer can connect $c$ separate
groups. Collect and print those representative pairs.

\begin{ideabox}[c Components Need Exactly $c-1$ Links]
Each new road can reduce the number of components by at most one, so going
from $c$ components to $1$ requires at least $c-1$ roads; connecting one
representative per component to a fixed anchor achieves that bound.
\end{ideabox}

\subsection*{Algorithm}

\begin{enumerate}
  \item Find components; record one representative node per component.
  \item Output $(\text{count}-1)$ as the number of roads.
  \item For each representative after the first, print a road from the
        first representative to it.
\end{enumerate}

\subsection*{Dry run}

\dryrun{$n=4$, roads $1\!-\!2$; cities $3$ and $4$ isolated.}

\begin{itemize}
  \item Components: $\{1,2\}$, $\{3\}$, $\{4\}$ -- reps $1,3,4$.
  \item Build $1\!-\!3$ and $1\!-\!4$: two roads. Answer $2$.
\end{itemize}

\subsection*{C++ implementation}

\begin{lstlisting}
#include <bits/stdc++.h>
using namespace std;

int main() {
    int n, m;
    cin >> n >> m;
    vector<vector<int>> adj(n + 1);
    for (int i = 0; i < m; i++) {
        int a, b; cin >> a >> b;
        adj[a].push_back(b);
        adj[b].push_back(a);
    }

    vector<bool> vis(n + 1, false);
    vector<int> reps;
    for (int s = 1; s <= n; s++) {
        if (vis[s]) continue;
        reps.push_back(s);                       // one representative per component
        queue<int> q; q.push(s); vis[s] = true;
        while (!q.empty()) {
            int u = q.front(); q.pop();
            for (int v : adj[u]) if (!vis[v]) { vis[v] = true; q.push(v); }
        }
    }

    cout << reps.size() - 1 << "\n";
    for (size_t i = 1; i < reps.size(); i++)
        cout << reps[0] << " " << reps[i] << "\n";
    return 0;
}
\end{lstlisting}

\begin{complexitybox}
\textbf{Time Complexity.} $O(n+m)$ -- one traversal to find components.
\textbf{Space Complexity.} $O(n+m)$.
\end{complexitybox}

\begin{takeawaybox}
Minimum edges to connect a graph is (components $-1$); connect one
representative per component to a common anchor. Recognising ``make it all
connected'' as component counting -- solvable by traversal or Disjoint Set
-- is the reusable idea.
\end{takeawaybox}
